% !PRINTED 184
\weberpage{184}

%<<CONT
und hieraus wird geschlossen, daß die Koeffizienten $A_\nu, B_\nu, C_\nu, D_\nu$
\gesperrt{ganze rationale Funktionen von $\varkappa^2$ vom Grade $\nu$ sind}, und
aus \S\,54, 2, 3, folgt überdies noch, daß $A_\nu, B_\nu$ bei geradem $\nu$
rational durch $\varkappa^2\varkappa'^2$ ausgedrückt sind, während sie bei ungeradem
$\nu$ gleich dem Produkt von $(\varkappa'^2-\varkappa^2)$ mit einer rationalen Funktion
 von $\varkappa^2\varkappa'^2$ sind.

Aus $B_\nu$ findet man $C_\nu$ mittels der Formel:

\weq{10}{C_\nu(\varkappa^2) = \varkappa'^{2\nu}B_\nu\left(\frac{-\varkappa^2}{\varkappa'^2}\right)}

und $D_\nu$ durch Anwendung von (6) und (8):

\weq{11}{D_\nu(\varkappa^2) = (-1)^\nu C_\nu(\varkappa'^2) = (-1)^\nu \varkappa^{2\nu}B_\nu\left(\frac{-\varkappa'^2}{\varkappa^2}\right).}

Da man aus (5) die Koeffizienten $A_\nu(0), B_\nu(0), C_\nu(0), D_\nu(0)$
leicht bestimmen kann, so ergeben sich aus dem hier entwickelten
Formelsystem die Koeffizienten $A_\nu, B_\nu, C_\nu, D_\nu$ ohne
weitere Rechnung bis $\nu=3$ einschließlich.

Man findet:
\[
\begin{aligned}
&A_0=1, \quad A_1=0, \qquad A_2=-\frac{2}{3}(1-\varkappa^2\varkappa'^2),\\
&B_0=1, \quad B_1=\frac{1}{3}(\varkappa'^2-\varkappa^2), \quad B_2=\frac{1}{3}(1+2\varkappa^2\varkappa'^2),\\
&C_0=1, \quad C_1=\frac{1}{3}(1+\varkappa^2), \quad C_2=\frac{1}{3}(\varkappa'^4-2\varkappa^2),\\
&D_0=1, \quad D_1=-\frac{1}{3}(1+\varkappa'^2), \quad D_2=\frac{1}{3}(\varkappa^4-2\varkappa'^2).\\
&A_3=-\frac{8}{9}(\varkappa'^2-\varkappa^2)(2+\varkappa^2\varkappa'^2),\\
&B_3=\frac{1}{9}(\varkappa'^2-\varkappa^2)(5-2\varkappa^2\varkappa'^2),\\
&C_3=\frac{1}{9}(1+\varkappa^2)(5\varkappa'^4+2\varkappa^2),\\
&D_3=-\frac{1}{9}(1+\varkappa'^2)(5\varkappa^4+2\varkappa'^2).
\end{aligned}
\]

Weitere Koeffizienten lassen sich auf demselben Wege, wenn
auch weitläufiger berechnen, indem man die Ausdrücke für die
$\sigma$-Funktionen in \S\,32 nach Potenzen von $q$ entwickelt und, wie
hier die niedrigste Potenz von $q$, so die höheren benutzt. \gesperrt{Weierstrass}
bedient sich zur rekurrenten Berechnung der Koeffizienten
gewisser partieller Differentialgleichungen, ähnlich der, die wir im
nächsten Paragraphen für die Funktion $\sigma$ ableiten werden.
